% ============================================================
% NEW SECTION for the NEW Chapter 2 (Related Work)
% Placement: second section, after perspectives_on_novelty.
%   % ============================================================
% NEW SECTION for the NEW Chapter 2 (Related Work)
% Placement: second section, after perspectives_on_novelty.
%   % ============================================================
% NEW SECTION for the NEW Chapter 2 (Related Work)
% Placement: second section, after perspectives_on_novelty.
%   % ============================================================
% NEW SECTION for the NEW Chapter 2 (Related Work)
% Placement: second section, after perspectives_on_novelty.
%   \input{chapters/chapter2/computational_open_world}
%
% This is the CS/AI computational review you asked to add, with
% the concept-drift-vs-novelty distinction as its centerpiece
% (sec:drift_vs_novelty). It is written to be accurate and fair
% to the continual-learning / concept-drift communities rather
% than to strawman them.
%
% Scope note: the hybrid planning-and-learning material here is
% deliberately high level and forward-points to Chapter 5's
% detailed related work, to avoid duplication. The open-world
% autonomy sections (sensing / planning-learning / execution-
% recovery) that follow in this chapter pick up the sequential-
% decision-making thread in depth.
%
% All citations verified or already present in bibliography.bib.
% New keys in bib_additions_v2.bib. VERIFY finn2017maml details.
% ============================================================

\section{Computational Approaches to Learning Under Change}
\label{sec:computational_open_world}

The previous section observed that adaptive systems, across disciplines, respond to unexpected experience in two modes: absorbing it within existing structure, or restructuring that model when absorption fails. This section reviews how artificial intelligence has formalized learning under change, and argues that the same divide separates the bulk of existing work from the problem this dissertation targets. We proceed from the standard learning setting and its stationarity assumptions, through the families of methods designed to relax those assumptions---continual learning, transfer, and concept-drift adaptation---to the distinction between distributional change and structural novelty that motivates the rest of the dissertation. We close with the exploratory drives that let agents seek out change, and a brief note on hybrid planning-and-learning, whose detailed treatment is deferred to \cref{ch:rapid_learn}.

\subsection{Reinforcement Learning and the Stationarity Assumption}
\label{sec:rl_background}

Reinforcement learning (RL) provides a general framework for an agent that learns to act by interacting with its environment, typically modeled as a Markov decision process (MDP) with fixed state and action spaces, transition dynamics, and reward function. Deep RL has produced strong results in games, control, and simulated robotics. These methods are powerful precisely because they exploit regularity: the policy or value function is optimized against a stationary environment, and performance guarantees rest on the assumption that the dynamics encountered at deployment match those seen during training. When that assumption holds, learned behavior can be highly effective; when it fails, a policy optimized for one environment can degrade sharply, because nothing in the standard formulation anticipates that the MDP itself might change. The remainder of this section concerns the methods that relax this assumption, and the extent to which they do so.

\subsection{Continual Learning, Transfer, and Concept Drift}
\label{sec:continual_transfer_drift}

Three overlapping research programs address learning when conditions are not fixed.

\emph{Continual and lifelong learning} were articulated as research agendas well before the deep learning era, motivated by agents that must accumulate and reuse knowledge over long deployments~\citep{ring1994continual, thrun1995lifelong}. The deep learning renaissance sharpened the central obstacle: neural networks trained on a sequence of tasks suffer \emph{catastrophic forgetting}, overwriting earlier competencies as new ones are acquired~\citep{parisi2019continual, hadsell2020embracing}. A large body of methods mitigates this, for example by protecting parameters important to previous tasks~\citep{kirkpatrick2017overcoming}. In the sequential-decision setting, continual reinforcement learning extends these concerns to non-stationary MDPs and streams of tasks~\citep{khetarpal2022continual}.

\emph{Transfer learning} asks how knowledge acquired in one task or environment can accelerate learning in another~\citep{taylor2009transfer}. \emph{Meta-learning} and meta-RL push this further, learning \emph{how to adapt} so that a new task can be solved from little experience---but typically under the assumption that the new task is drawn from the same distribution as those seen during meta-training~\citep{finn2017maml, yu2020metaworld}.

\emph{Concept-drift adaptation}, studied largely in the data-stream and supervised-learning communities, addresses the case where the distribution generating an agent's data changes over time~\citep{widmer1996learning}. The problem is usually formalized as a change in the joint distribution $P_t(X, Y)$ with time $t$, and a well-developed taxonomy distinguishes, for example, \emph{real drift} (a change in $P(Y \mid X)$) from \emph{virtual drift} (a change in $P(X)$ alone), and sudden from gradual and recurring drift~\citep{gama2014survey}. Methods pair a drift-detection mechanism with an adaptation strategy such as windowing, instance weighting, or ensemble reconfiguration~\citep{lu2019concept}.

These programs differ in emphasis, but they share a common structural assumption, which the next subsection makes explicit.

\subsection{Distributional Change Versus Structural Novelty}
\label{sec:drift_vs_novelty}

Continual learning, transfer, and concept-drift adaptation are all sometimes described as forms of ``novelty handling,'' and each is valuable within its scope. It is worth being precise, however, about the kind of change each assumes, because it differs from the kind that open-world environments impose.

In concept drift, what changes is the \emph{distribution} over a fixed feature space $X$ and label space $Y$: the variables that describe the world, and the set of possible outcomes, remain the same, while their statistics move~\citep{gama2014survey, lu2019concept}. In continual and meta-reinforcement learning, the state and action spaces are likewise typically held fixed across the stream of tasks; what varies is the reward, the dynamics, or the task identity, and adaptation proceeds by re-estimating or re-weighting a policy defined over that fixed representation~\citep{khetarpal2022continual, finn2017maml}. In all of these settings, the agent's \emph{representational vocabulary} is given in advance, and learning under change means adjusting parameters within it. This is, in the terms of \cref{sec:novelty_psychology}, an assimilation regime: new experience is accommodated by updating a model whose structure is already adequate to describe it.

Open-world novelty need not respect this assumption. A previously unseen object type may appear, adding a category the agent has no variable for; an action may acquire a new precondition, changing the transition structure qualitatively rather than shifting a probability; the set of entities, relations, or operators relevant to the task may expand. Changes of this kind alter the support and structure of the agent's representation, not merely the distribution defined over it. Recognizing that the world has changed is not enough; the agent must \emph{extend} its representation---introduce a new symbol, entity, or skill---to describe and act on the change at all. In the vocabulary of \cref{sec:novelty_psychology}, this is accommodation rather than assimilation, and it is not what methods built for distributional change are designed to do. This is why open-world novelty is not simply ``more non-stationarity'': quantitatively larger drift within a fixed representation remains a different problem from a change that the representation cannot express.

The boundary is not absolute, and parts of the field have moved toward it. Open-set and open-world recognition study classifiers that must reject or incorporate categories absent from training~\citep{boult2021towards}, and class-incremental learning studies models that add new classes over time. These lines confront structural change directly, but predominantly in a perception-and-classification setting. What remains comparatively underexplored---and what this dissertation targets---is structural novelty in the full sequential-decision-making setting, where a change can invalidate not only a category boundary but the executability of a plan, and where recovery requires acquiring new behavior rather than only re-labeling an input. The benchmarks of \cref{ch:benchmarks} are designed to inject exactly such structural changes under controlled conditions, and the adaptation methods of \cref{ch:rapid_learn} are designed to extend an agent's representation in response to them.

\subsection{Intrinsic Motivation and Exploration}
\label{sec:intrinsic_motivation}

A complementary line of work concerns not how an agent responds to change, but what drives it to seek change out. Intrinsic-motivation and artificial-curiosity methods give agents an internal incentive to explore, formalizing the psychological observations of \cref{sec:novelty_psychology} as computational objectives: agents are rewarded for prediction error, learning progress, or information gain~\citep{schmidhuber1991curious, oudeyer2007intrinsic, pathak2017curiosity}. These mechanisms are largely orthogonal to the present work but naturally complementary to it: this dissertation is concerned mainly with what an agent should do once change has manifested as task-relevant failure, whereas curiosity-style objectives are well suited to driving the exploration phases through which new behavior is discovered.

\subsection{Hybrid Planning and Learning}
\label{sec:hybrid_planning_learning_rw}

Finally, a substantial body of work integrates symbolic planning with learned control, combining the compositional structure and generalization of planning with the flexibility of learning. Because this material bears directly on the adaptation method developed later in the dissertation, its detailed treatment---including task-and-motion planning, learned symbolic operators, plan repair, and prior open-world architectures---is deferred to the related-work discussion of \cref{ch:rapid_learn}. For the purposes of this chapter, the relevant point is that structured, symbolic representations provide a natural mechanism for \emph{localizing} change: when behavior is organized into named units with explicit preconditions and effects, a failure can be attributed to a specific component rather than diffused across a monolithic policy. This property is central to the methods of \cref{ch:rapid_learn} and connects back to the representational theme running through this chapter.

%
% This is the CS/AI computational review you asked to add, with
% the concept-drift-vs-novelty distinction as its centerpiece
% (sec:drift_vs_novelty). It is written to be accurate and fair
% to the continual-learning / concept-drift communities rather
% than to strawman them.
%
% Scope note: the hybrid planning-and-learning material here is
% deliberately high level and forward-points to Chapter 5's
% detailed related work, to avoid duplication. The open-world
% autonomy sections (sensing / planning-learning / execution-
% recovery) that follow in this chapter pick up the sequential-
% decision-making thread in depth.
%
% All citations verified or already present in bibliography.bib.
% New keys in bib_additions_v2.bib. VERIFY finn2017maml details.
% ============================================================

\section{Computational Approaches to Learning Under Change}
\label{sec:computational_open_world}

The previous section observed that adaptive systems, across disciplines, respond to unexpected experience in two modes: absorbing it within existing structure, or restructuring that model when absorption fails. This section reviews how artificial intelligence has formalized learning under change, and argues that the same divide separates the bulk of existing work from the problem this dissertation targets. We proceed from the standard learning setting and its stationarity assumptions, through the families of methods designed to relax those assumptions---continual learning, transfer, and concept-drift adaptation---to the distinction between distributional change and structural novelty that motivates the rest of the dissertation. We close with the exploratory drives that let agents seek out change, and a brief note on hybrid planning-and-learning, whose detailed treatment is deferred to \cref{ch:rapid_learn}.

\subsection{Reinforcement Learning and the Stationarity Assumption}
\label{sec:rl_background}

Reinforcement learning (RL) provides a general framework for an agent that learns to act by interacting with its environment, typically modeled as a Markov decision process (MDP) with fixed state and action spaces, transition dynamics, and reward function. Deep RL has produced strong results in games, control, and simulated robotics. These methods are powerful precisely because they exploit regularity: the policy or value function is optimized against a stationary environment, and performance guarantees rest on the assumption that the dynamics encountered at deployment match those seen during training. When that assumption holds, learned behavior can be highly effective; when it fails, a policy optimized for one environment can degrade sharply, because nothing in the standard formulation anticipates that the MDP itself might change. The remainder of this section concerns the methods that relax this assumption, and the extent to which they do so.

\subsection{Continual Learning, Transfer, and Concept Drift}
\label{sec:continual_transfer_drift}

Three overlapping research programs address learning when conditions are not fixed.

\emph{Continual and lifelong learning} were articulated as research agendas well before the deep learning era, motivated by agents that must accumulate and reuse knowledge over long deployments~\citep{ring1994continual, thrun1995lifelong}. The deep learning renaissance sharpened the central obstacle: neural networks trained on a sequence of tasks suffer \emph{catastrophic forgetting}, overwriting earlier competencies as new ones are acquired~\citep{parisi2019continual, hadsell2020embracing}. A large body of methods mitigates this, for example by protecting parameters important to previous tasks~\citep{kirkpatrick2017overcoming}. In the sequential-decision setting, continual reinforcement learning extends these concerns to non-stationary MDPs and streams of tasks~\citep{khetarpal2022continual}.

\emph{Transfer learning} asks how knowledge acquired in one task or environment can accelerate learning in another~\citep{taylor2009transfer}. \emph{Meta-learning} and meta-RL push this further, learning \emph{how to adapt} so that a new task can be solved from little experience---but typically under the assumption that the new task is drawn from the same distribution as those seen during meta-training~\citep{finn2017maml, yu2020metaworld}.

\emph{Concept-drift adaptation}, studied largely in the data-stream and supervised-learning communities, addresses the case where the distribution generating an agent's data changes over time~\citep{widmer1996learning}. The problem is usually formalized as a change in the joint distribution $P_t(X, Y)$ with time $t$, and a well-developed taxonomy distinguishes, for example, \emph{real drift} (a change in $P(Y \mid X)$) from \emph{virtual drift} (a change in $P(X)$ alone), and sudden from gradual and recurring drift~\citep{gama2014survey}. Methods pair a drift-detection mechanism with an adaptation strategy such as windowing, instance weighting, or ensemble reconfiguration~\citep{lu2019concept}.

These programs differ in emphasis, but they share a common structural assumption, which the next subsection makes explicit.

\subsection{Distributional Change Versus Structural Novelty}
\label{sec:drift_vs_novelty}

Continual learning, transfer, and concept-drift adaptation are all sometimes described as forms of ``novelty handling,'' and each is valuable within its scope. It is worth being precise, however, about the kind of change each assumes, because it differs from the kind that open-world environments impose.

In concept drift, what changes is the \emph{distribution} over a fixed feature space $X$ and label space $Y$: the variables that describe the world, and the set of possible outcomes, remain the same, while their statistics move~\citep{gama2014survey, lu2019concept}. In continual and meta-reinforcement learning, the state and action spaces are likewise typically held fixed across the stream of tasks; what varies is the reward, the dynamics, or the task identity, and adaptation proceeds by re-estimating or re-weighting a policy defined over that fixed representation~\citep{khetarpal2022continual, finn2017maml}. In all of these settings, the agent's \emph{representational vocabulary} is given in advance, and learning under change means adjusting parameters within it. This is, in the terms of \cref{sec:novelty_psychology}, an assimilation regime: new experience is accommodated by updating a model whose structure is already adequate to describe it.

Open-world novelty need not respect this assumption. A previously unseen object type may appear, adding a category the agent has no variable for; an action may acquire a new precondition, changing the transition structure qualitatively rather than shifting a probability; the set of entities, relations, or operators relevant to the task may expand. Changes of this kind alter the support and structure of the agent's representation, not merely the distribution defined over it. Recognizing that the world has changed is not enough; the agent must \emph{extend} its representation---introduce a new symbol, entity, or skill---to describe and act on the change at all. In the vocabulary of \cref{sec:novelty_psychology}, this is accommodation rather than assimilation, and it is not what methods built for distributional change are designed to do. This is why open-world novelty is not simply ``more non-stationarity'': quantitatively larger drift within a fixed representation remains a different problem from a change that the representation cannot express.

The boundary is not absolute, and parts of the field have moved toward it. Open-set and open-world recognition study classifiers that must reject or incorporate categories absent from training~\citep{boult2021towards}, and class-incremental learning studies models that add new classes over time. These lines confront structural change directly, but predominantly in a perception-and-classification setting. What remains comparatively underexplored---and what this dissertation targets---is structural novelty in the full sequential-decision-making setting, where a change can invalidate not only a category boundary but the executability of a plan, and where recovery requires acquiring new behavior rather than only re-labeling an input. The benchmarks of \cref{ch:benchmarks} are designed to inject exactly such structural changes under controlled conditions, and the adaptation methods of \cref{ch:rapid_learn} are designed to extend an agent's representation in response to them.

\subsection{Intrinsic Motivation and Exploration}
\label{sec:intrinsic_motivation}

A complementary line of work concerns not how an agent responds to change, but what drives it to seek change out. Intrinsic-motivation and artificial-curiosity methods give agents an internal incentive to explore, formalizing the psychological observations of \cref{sec:novelty_psychology} as computational objectives: agents are rewarded for prediction error, learning progress, or information gain~\citep{schmidhuber1991curious, oudeyer2007intrinsic, pathak2017curiosity}. These mechanisms are largely orthogonal to the present work but naturally complementary to it: this dissertation is concerned mainly with what an agent should do once change has manifested as task-relevant failure, whereas curiosity-style objectives are well suited to driving the exploration phases through which new behavior is discovered.

\subsection{Hybrid Planning and Learning}
\label{sec:hybrid_planning_learning_rw}

Finally, a substantial body of work integrates symbolic planning with learned control, combining the compositional structure and generalization of planning with the flexibility of learning. Because this material bears directly on the adaptation method developed later in the dissertation, its detailed treatment---including task-and-motion planning, learned symbolic operators, plan repair, and prior open-world architectures---is deferred to the related-work discussion of \cref{ch:rapid_learn}. For the purposes of this chapter, the relevant point is that structured, symbolic representations provide a natural mechanism for \emph{localizing} change: when behavior is organized into named units with explicit preconditions and effects, a failure can be attributed to a specific component rather than diffused across a monolithic policy. This property is central to the methods of \cref{ch:rapid_learn} and connects back to the representational theme running through this chapter.

%
% This is the CS/AI computational review you asked to add, with
% the concept-drift-vs-novelty distinction as its centerpiece
% (sec:drift_vs_novelty). It is written to be accurate and fair
% to the continual-learning / concept-drift communities rather
% than to strawman them.
%
% Scope note: the hybrid planning-and-learning material here is
% deliberately high level and forward-points to Chapter 5's
% detailed related work, to avoid duplication. The open-world
% autonomy sections (sensing / planning-learning / execution-
% recovery) that follow in this chapter pick up the sequential-
% decision-making thread in depth.
%
% All citations verified or already present in bibliography.bib.
% New keys in bib_additions_v2.bib. VERIFY finn2017maml details.
% ============================================================

\section{Computational Approaches to Learning Under Change}
\label{sec:computational_open_world}

The previous section observed that adaptive systems, across disciplines, respond to unexpected experience in two modes: absorbing it within existing structure, or restructuring that model when absorption fails. This section reviews how artificial intelligence has formalized learning under change, and argues that the same divide separates the bulk of existing work from the problem this dissertation targets. We proceed from the standard learning setting and its stationarity assumptions, through the families of methods designed to relax those assumptions---continual learning, transfer, and concept-drift adaptation---to the distinction between distributional change and structural novelty that motivates the rest of the dissertation. We close with the exploratory drives that let agents seek out change, and a brief note on hybrid planning-and-learning, whose detailed treatment is deferred to \cref{ch:rapid_learn}.

\subsection{Reinforcement Learning and the Stationarity Assumption}
\label{sec:rl_background}

Reinforcement learning (RL) provides a general framework for an agent that learns to act by interacting with its environment, typically modeled as a Markov decision process (MDP) with fixed state and action spaces, transition dynamics, and reward function. Deep RL has produced strong results in games, control, and simulated robotics. These methods are powerful precisely because they exploit regularity: the policy or value function is optimized against a stationary environment, and performance guarantees rest on the assumption that the dynamics encountered at deployment match those seen during training. When that assumption holds, learned behavior can be highly effective; when it fails, a policy optimized for one environment can degrade sharply, because nothing in the standard formulation anticipates that the MDP itself might change. The remainder of this section concerns the methods that relax this assumption, and the extent to which they do so.

\subsection{Continual Learning, Transfer, and Concept Drift}
\label{sec:continual_transfer_drift}

Three overlapping research programs address learning when conditions are not fixed.

\emph{Continual and lifelong learning} were articulated as research agendas well before the deep learning era, motivated by agents that must accumulate and reuse knowledge over long deployments~\citep{ring1994continual, thrun1995lifelong}. The deep learning renaissance sharpened the central obstacle: neural networks trained on a sequence of tasks suffer \emph{catastrophic forgetting}, overwriting earlier competencies as new ones are acquired~\citep{parisi2019continual, hadsell2020embracing}. A large body of methods mitigates this, for example by protecting parameters important to previous tasks~\citep{kirkpatrick2017overcoming}. In the sequential-decision setting, continual reinforcement learning extends these concerns to non-stationary MDPs and streams of tasks~\citep{khetarpal2022continual}.

\emph{Transfer learning} asks how knowledge acquired in one task or environment can accelerate learning in another~\citep{taylor2009transfer}. \emph{Meta-learning} and meta-RL push this further, learning \emph{how to adapt} so that a new task can be solved from little experience---but typically under the assumption that the new task is drawn from the same distribution as those seen during meta-training~\citep{finn2017maml, yu2020metaworld}.

\emph{Concept-drift adaptation}, studied largely in the data-stream and supervised-learning communities, addresses the case where the distribution generating an agent's data changes over time~\citep{widmer1996learning}. The problem is usually formalized as a change in the joint distribution $P_t(X, Y)$ with time $t$, and a well-developed taxonomy distinguishes, for example, \emph{real drift} (a change in $P(Y \mid X)$) from \emph{virtual drift} (a change in $P(X)$ alone), and sudden from gradual and recurring drift~\citep{gama2014survey}. Methods pair a drift-detection mechanism with an adaptation strategy such as windowing, instance weighting, or ensemble reconfiguration~\citep{lu2019concept}.

These programs differ in emphasis, but they share a common structural assumption, which the next subsection makes explicit.

\subsection{Distributional Change Versus Structural Novelty}
\label{sec:drift_vs_novelty}

Continual learning, transfer, and concept-drift adaptation are all sometimes described as forms of ``novelty handling,'' and each is valuable within its scope. It is worth being precise, however, about the kind of change each assumes, because it differs from the kind that open-world environments impose.

In concept drift, what changes is the \emph{distribution} over a fixed feature space $X$ and label space $Y$: the variables that describe the world, and the set of possible outcomes, remain the same, while their statistics move~\citep{gama2014survey, lu2019concept}. In continual and meta-reinforcement learning, the state and action spaces are likewise typically held fixed across the stream of tasks; what varies is the reward, the dynamics, or the task identity, and adaptation proceeds by re-estimating or re-weighting a policy defined over that fixed representation~\citep{khetarpal2022continual, finn2017maml}. In all of these settings, the agent's \emph{representational vocabulary} is given in advance, and learning under change means adjusting parameters within it. This is, in the terms of \cref{sec:novelty_psychology}, an assimilation regime: new experience is accommodated by updating a model whose structure is already adequate to describe it.

Open-world novelty need not respect this assumption. A previously unseen object type may appear, adding a category the agent has no variable for; an action may acquire a new precondition, changing the transition structure qualitatively rather than shifting a probability; the set of entities, relations, or operators relevant to the task may expand. Changes of this kind alter the support and structure of the agent's representation, not merely the distribution defined over it. Recognizing that the world has changed is not enough; the agent must \emph{extend} its representation---introduce a new symbol, entity, or skill---to describe and act on the change at all. In the vocabulary of \cref{sec:novelty_psychology}, this is accommodation rather than assimilation, and it is not what methods built for distributional change are designed to do. This is why open-world novelty is not simply ``more non-stationarity'': quantitatively larger drift within a fixed representation remains a different problem from a change that the representation cannot express.

The boundary is not absolute, and parts of the field have moved toward it. Open-set and open-world recognition study classifiers that must reject or incorporate categories absent from training~\citep{boult2021towards}, and class-incremental learning studies models that add new classes over time. These lines confront structural change directly, but predominantly in a perception-and-classification setting. What remains comparatively underexplored---and what this dissertation targets---is structural novelty in the full sequential-decision-making setting, where a change can invalidate not only a category boundary but the executability of a plan, and where recovery requires acquiring new behavior rather than only re-labeling an input. The benchmarks of \cref{ch:benchmarks} are designed to inject exactly such structural changes under controlled conditions, and the adaptation methods of \cref{ch:rapid_learn} are designed to extend an agent's representation in response to them.

\subsection{Intrinsic Motivation and Exploration}
\label{sec:intrinsic_motivation}

A complementary line of work concerns not how an agent responds to change, but what drives it to seek change out. Intrinsic-motivation and artificial-curiosity methods give agents an internal incentive to explore, formalizing the psychological observations of \cref{sec:novelty_psychology} as computational objectives: agents are rewarded for prediction error, learning progress, or information gain~\citep{schmidhuber1991curious, oudeyer2007intrinsic, pathak2017curiosity}. These mechanisms are largely orthogonal to the present work but naturally complementary to it: this dissertation is concerned mainly with what an agent should do once change has manifested as task-relevant failure, whereas curiosity-style objectives are well suited to driving the exploration phases through which new behavior is discovered.

\subsection{Hybrid Planning and Learning}
\label{sec:hybrid_planning_learning_rw}

Finally, a substantial body of work integrates symbolic planning with learned control, combining the compositional structure and generalization of planning with the flexibility of learning. Because this material bears directly on the adaptation method developed later in the dissertation, its detailed treatment---including task-and-motion planning, learned symbolic operators, plan repair, and prior open-world architectures---is deferred to the related-work discussion of \cref{ch:rapid_learn}. For the purposes of this chapter, the relevant point is that structured, symbolic representations provide a natural mechanism for \emph{localizing} change: when behavior is organized into named units with explicit preconditions and effects, a failure can be attributed to a specific component rather than diffused across a monolithic policy. This property is central to the methods of \cref{ch:rapid_learn} and connects back to the representational theme running through this chapter.

%
% This is the CS/AI computational review you asked to add, with
% the concept-drift-vs-novelty distinction as its centerpiece
% (sec:drift_vs_novelty). It is written to be accurate and fair
% to the continual-learning / concept-drift communities rather
% than to strawman them.
%
% Scope note: the hybrid planning-and-learning material here is
% deliberately high level and forward-points to Chapter 5's
% detailed related work, to avoid duplication. The open-world
% autonomy sections (sensing / planning-learning / execution-
% recovery) that follow in this chapter pick up the sequential-
% decision-making thread in depth.
%
% All citations verified or already present in bibliography.bib.
% New keys in bib_additions_v2.bib. VERIFY finn2017maml details.
% ============================================================

\section{Computational Approaches to Learning Under Change}
\label{sec:computational_open_world}

The previous section observed that adaptive systems, across disciplines, respond to unexpected experience in two modes: absorbing it within existing structure, or restructuring that model when absorption fails. This section reviews how artificial intelligence has formalized learning under change, and argues that the same divide separates the bulk of existing work from the problem this dissertation targets. We proceed from the standard learning setting and its stationarity assumptions, through the families of methods designed to relax those assumptions---continual learning, transfer, and concept-drift adaptation---to the distinction between distributional change and structural novelty that motivates the rest of the dissertation. We close with the exploratory drives that let agents seek out change, and a brief note on hybrid planning-and-learning, whose detailed treatment is deferred to \cref{ch:rapid_learn}.

\subsection{Reinforcement Learning and the Stationarity Assumption}
\label{sec:rl_background}

Reinforcement learning (RL) provides a general framework for an agent that learns to act by interacting with its environment, typically modeled as a Markov decision process (MDP) with fixed state and action spaces, transition dynamics, and reward function. Deep RL has produced strong results in games, control, and simulated robotics. These methods are powerful precisely because they exploit regularity: the policy or value function is optimized against a stationary environment, and performance guarantees rest on the assumption that the dynamics encountered at deployment match those seen during training. When that assumption holds, learned behavior can be highly effective; when it fails, a policy optimized for one environment can degrade sharply, because nothing in the standard formulation anticipates that the MDP itself might change. The remainder of this section concerns the methods that relax this assumption, and the extent to which they do so.

\subsection{Continual Learning, Transfer, and Concept Drift}
\label{sec:continual_transfer_drift}

Three overlapping research programs address learning when conditions are not fixed.

\emph{Continual and lifelong learning} were articulated as research agendas well before the deep learning era, motivated by agents that must accumulate and reuse knowledge over long deployments~\citep{ring1994continual, thrun1995lifelong}. The deep learning renaissance sharpened the central obstacle: neural networks trained on a sequence of tasks suffer \emph{catastrophic forgetting}, overwriting earlier competencies as new ones are acquired~\citep{parisi2019continual, hadsell2020embracing}. A large body of methods mitigates this, for example by protecting parameters important to previous tasks~\citep{kirkpatrick2017overcoming}. In the sequential-decision setting, continual reinforcement learning extends these concerns to non-stationary MDPs and streams of tasks~\citep{khetarpal2022continual}.

\emph{Transfer learning} asks how knowledge acquired in one task or environment can accelerate learning in another~\citep{taylor2009transfer}. \emph{Meta-learning} and meta-RL push this further, learning \emph{how to adapt} so that a new task can be solved from little experience---but typically under the assumption that the new task is drawn from the same distribution as those seen during meta-training~\citep{finn2017maml, yu2020metaworld}.

\emph{Concept-drift adaptation}, studied largely in the data-stream and supervised-learning communities, addresses the case where the distribution generating an agent's data changes over time~\citep{widmer1996learning}. The problem is usually formalized as a change in the joint distribution $P_t(X, Y)$ with time $t$, and a well-developed taxonomy distinguishes, for example, \emph{real drift} (a change in $P(Y \mid X)$) from \emph{virtual drift} (a change in $P(X)$ alone), and sudden from gradual and recurring drift~\citep{gama2014survey}. Methods pair a drift-detection mechanism with an adaptation strategy such as windowing, instance weighting, or ensemble reconfiguration~\citep{lu2019concept}.

These programs differ in emphasis, but they share a common structural assumption, which the next subsection makes explicit.

\subsection{Distributional Change Versus Structural Novelty}
\label{sec:drift_vs_novelty}

Continual learning, transfer, and concept-drift adaptation are all sometimes described as forms of ``novelty handling,'' and each is valuable within its scope. It is worth being precise, however, about the kind of change each assumes, because it differs from the kind that open-world environments impose.

In concept drift, what changes is the \emph{distribution} over a fixed feature space $X$ and label space $Y$: the variables that describe the world, and the set of possible outcomes, remain the same, while their statistics move~\citep{gama2014survey, lu2019concept}. In continual and meta-reinforcement learning, the state and action spaces are likewise typically held fixed across the stream of tasks; what varies is the reward, the dynamics, or the task identity, and adaptation proceeds by re-estimating or re-weighting a policy defined over that fixed representation~\citep{khetarpal2022continual, finn2017maml}. In all of these settings, the agent's \emph{representational vocabulary} is given in advance, and learning under change means adjusting parameters within it. This is, in the terms of \cref{sec:novelty_psychology}, an assimilation regime: new experience is accommodated by updating a model whose structure is already adequate to describe it.

Open-world novelty need not respect this assumption. A previously unseen object type may appear, adding a category the agent has no variable for; an action may acquire a new precondition, changing the transition structure qualitatively rather than shifting a probability; the set of entities, relations, or operators relevant to the task may expand. Changes of this kind alter the support and structure of the agent's representation, not merely the distribution defined over it. Recognizing that the world has changed is not enough; the agent must \emph{extend} its representation---introduce a new symbol, entity, or skill---to describe and act on the change at all. In the vocabulary of \cref{sec:novelty_psychology}, this is accommodation rather than assimilation, and it is not what methods built for distributional change are designed to do. This is why open-world novelty is not simply ``more non-stationarity'': quantitatively larger drift within a fixed representation remains a different problem from a change that the representation cannot express.

The boundary is not absolute, and parts of the field have moved toward it. Open-set and open-world recognition study classifiers that must reject or incorporate categories absent from training~\citep{boult2021towards}, and class-incremental learning studies models that add new classes over time. These lines confront structural change directly, but predominantly in a perception-and-classification setting. What remains comparatively underexplored---and what this dissertation targets---is structural novelty in the full sequential-decision-making setting, where a change can invalidate not only a category boundary but the executability of a plan, and where recovery requires acquiring new behavior rather than only re-labeling an input. The benchmarks of \cref{ch:benchmarks} are designed to inject exactly such structural changes under controlled conditions, and the adaptation methods of \cref{ch:rapid_learn} are designed to extend an agent's representation in response to them.

\subsection{Intrinsic Motivation and Exploration}
\label{sec:intrinsic_motivation}

A complementary line of work concerns not how an agent responds to change, but what drives it to seek change out. Intrinsic-motivation and artificial-curiosity methods give agents an internal incentive to explore, formalizing the psychological observations of \cref{sec:novelty_psychology} as computational objectives: agents are rewarded for prediction error, learning progress, or information gain~\citep{schmidhuber1991curious, oudeyer2007intrinsic, pathak2017curiosity}. These mechanisms are largely orthogonal to the present work but naturally complementary to it: this dissertation is concerned mainly with what an agent should do once change has manifested as task-relevant failure, whereas curiosity-style objectives are well suited to driving the exploration phases through which new behavior is discovered.

\subsection{Hybrid Planning and Learning}
\label{sec:hybrid_planning_learning_rw}

Finally, a substantial body of work integrates symbolic planning with learned control, combining the compositional structure and generalization of planning with the flexibility of learning. Because this material bears directly on the adaptation method developed later in the dissertation, its detailed treatment---including task-and-motion planning, learned symbolic operators, plan repair, and prior open-world architectures---is deferred to the related-work discussion of \cref{ch:rapid_learn}. For the purposes of this chapter, the relevant point is that structured, symbolic representations provide a natural mechanism for \emph{localizing} change: when behavior is organized into named units with explicit preconditions and effects, a failure can be attributed to a specific component rather than diffused across a monolithic policy. This property is central to the methods of \cref{ch:rapid_learn} and connects back to the representational theme running through this chapter.
